% No \label here. \exercises is a starred section, so a label attached to it
% takes the number of the last *numbered* section instead - it silently
% resolved to 9.6 in a first draft, which is the closing section, and the build
% reported nothing because the label was defined, just wrong. Prose refers to
% these by name.
\exercises

\tierunderstand

\begin{bt}
  \item Dẫn \eqref{eq:ch09-phan-attention} lại từ đầu, viết rõ chi phí mỗi
    token mỗi tầng của bốn phép chiếu, của mạng truyền thẳng, và của hai tích
    attention. Rồi trả lời câu đắt hơn: vì sao số tầng biến mất khỏi kết quả,
    và điều đó còn đúng không nếu một mô hình đặt \(\dff\) khác nhau ở các tầng
    khác nhau?

  \item Mục~\ref{sec:ch09-sau-nd} đo được \(6ND\) sai 2.8\% ở GPT-3 và 11\% ở
    GPT-1, dù GPT-3 lớn hơn nghìn lần. Giải thích bằng
    \eqref{eq:ch09-phan-attention}. Rồi tính mà không chạy gì: một mô hình
    \(\dmodel = 4096\) phải lấy ngữ cảnh bằng bao nhiêu để \(6ND\) sai đúng
    10\%?

  \item Bảng 1 của Kaplan đếm một trong hai tích bậc hai. Chỉ ra hàng nào của
    bảng đáng lẽ phải có mặt và không có, và viết ra ô \code{FLOPs per Token}
    của nó. Rồi trả lời: nếu sửa bảng cho đủ, hàng \code{Total} đổi thế nào, và
    điều kiện \(d_{\mathrm{model}} \gg n_{\mathrm{ctx}}/12\) mà bài viết ở mục
    2.1 có đổi theo không?

  \item Từ \eqref{eq:ch09-chinchilla}, chứng minh rằng cực tiểu của \(L\) dưới
    ràng buộc \(6ND = C\) cho \(N_{\mathrm{opt}} \propto C^{a}\) với
    \(a = \beta/(\alpha+\beta)\). Rồi nói vì sao \(a + b = 1\) là kết luận của
    dạng hàm chứ không phải một phát hiện thực nghiệm.

  \item Mục~\ref{sec:ch09-phat-kien} đo được masked language model thu 76 đích
    trên một chuỗi 512 token còn \tn{mô hình ngôn ngữ}{language model} thu 512.
    Nhưng hai mục tiêu
    ấy không cho ra hai mô hình so sánh trực tiếp được. Nêu ít nhất hai lý do,
    và với mỗi lý do nói phép đo nào sẽ khử được nó.

  \item Bảng 5 mục 5.1 của bài BERT cho SQuAD F1 là 88.5 ở cấu hình đầy đủ,
    87.9 khi bỏ next-sentence prediction, và 77.8 khi bỏ cả hai chiều lẫn
    next-sentence prediction. Tách hai đóng góp ra từ ba con số ấy, rồi nói
    phép tách của bạn dựa trên giả thiết nào và giả thiết ấy có được bảng kiểm
    hay không.

  \item Kaplan và Hoffmann định nghĩa \(N\) ngược nhau, một bên trừ embedding
    và một bên cộng. Dùng bảng ở mục~\ref{sec:ch09-sau-nd}: ở GPT-1 và ở GPT-3,
    một quy luật khớp theo định nghĩa này mà áp theo định nghĩa kia sẽ lệch bao
    nhiêu? Rồi trả lời câu thật sự khó: chỗ lệch ấy có giải thích được khoảng
    cách 0.73 với 0.50 không, và bạn kiểm bằng cách nào?

  \item Mục~\ref{sec:ch09-cung-mot-khoi} giải ngược bảng 2 của GPT-2 ra bốn
    kích thước từ điển quanh 40,500. Phép giải ấy giả thiết bề rộng, độ sâu và
    \(\dff\) đều đúng như bảng in. Nêu một \tn{giả thuyết}{hypothesis} khác
    giải thích được cùng
    bốn khoảng cách ấy, rồi thiết kế một phép kiểm tách hai giả thuyết bằng
    đúng những con số bài in.
\end{bt}

\tierapply

Trạng thái xuất phát cho cả năm bài: clone repo companion, đứng ở thư mục gốc
của nó, và checkout đúng tag của chương này.

\begin{minted}{console}
$ git clone https://github.com/Giang-Dang/rnn-to-transformer-lab.git
$ cd rnn-to-transformer-lab
$ git checkout ch09
$ conda env create -f environment.yml
$ conda activate rnn-to-transformer-lab
$ python verify.py --only ch09
\end{minted}

Dòng cuối phải in ra \code{verify: ok}. Ba mục của chương này chạy hết dưới năm
giây cộng lại, vì không mục nào huấn luyện mô hình; nếu máy bạn mất lâu hơn
nhiều thì gần như chắc chắn là thời gian khởi động của trình thông dịch chứ
không phải phép tính.

Vẫn cảnh báo cũ: đừng gọi các script qua \code{conda run}.

\begin{bt}
  \item Dựng lại bảng đầu tiên của chương trên một mô hình bạn tra được siêu
    tham số. Thêm một hàm vào \code{scaling.py} cho kiến trúc ấy nếu nó khác
    hình dạng khối ở đây, rồi so số bạn dựng với số nó công bố. Nói rõ bạn lấy
    siêu tham số từ đâu, và nếu con số không khớp thì đưa ra ít nhất hai cách
    đọc trước khi kết luận cách nào đúng.

  \item Đo giới hạn của \code{block\_params}. Hàm ấy giả thiết mạng truyền
    thẳng có hai ma trận. Sửa nó cho kiến trúc dùng cổng, nơi tầng trong có ba
    ma trận thay vì hai, rồi chạy lại \code{ch09\_counts.py} và nói bảng nào
    đổi. Rồi chạy bộ test của chương:

    \begin{minted}{console}
$ python -m pytest -q tests/test_ch09.py
    \end{minted}

    Test nào bắt được thay đổi của bạn, test nào không, và vì sao?

  \item Kiểm \eqref{eq:ch09-phan-attention} bằng đồng hồ chứ không bằng số học.
    \code{ch08\_clock.py} của chương trước đo thời gian thật theo \(n\). Dùng
    nó để hỏi: phần thời gian attention chiếm có đi theo \(n/(6d+n)\) không?
    Biết trước là không, và mục~\ref{sec:ch08-cho-vo} đã nói vì sao. Chạy đi,
    rồi đo xem hai đường tách nhau ở đâu và nói chỗ tách ấy khớp với đại lượng
    nào của chương trước.

  \item Đẩy \eqref{eq:ch09-chinchilla} tới chỗ nó vô nghĩa. Hàm ấy đơn điệu
    giảm theo cả hai biến, nên nó luôn thích mô hình lớn hơn và dữ liệu nhiều
    hơn nếu không có ràng buộc. Viết một script quét \(N\) từ \(10^{6}\) tới
    \(10^{15}\) ở \(D\) cố định và tìm chỗ \(L\) rơi xuống dưới \(E = 1.69\).
    Chỗ ấy có tồn tại không? Trả lời câu tiếp: điều đó nói gì về việc dùng hàm
    này ngoài dải nó được khớp?

  \item Dựng lại bảng A4 của Hoffmann, biết trước rằng
    mục~\ref{sec:ch09-hai-quy-luat} đã thử và không được. Cài phép đếm của phụ
    lục F từng dòng một, rồi quét độ dài chuỗi và kích thước từ điển trên một
    lưới rộng hơn lưới chín điểm tôi thử. Nếu bạn tìm được cặp cho lại cả sáu
    con số, viết ra. Nếu không, hãy nói rõ hơn tôi đã nói: dải nào của hai tham
    số ấy bị loại trừ, và bị loại trừ bởi hàng nào của bảng?
\end{bt}

\tierextend

\begin{bt}
  \item Bài Hoffmann chỉ ra bài Kaplan sai vì một chi tiết của lịch learning
    rate, và chi tiết ấy chỉ lộ ra khi ai đó chạy lại phép khớp với lịch đặt
    đúng. Câu hỏi: còn chi tiết nào cùng loại trong chính bài Hoffmann? Chọn
    một quyết định phương pháp của bài, dựng một phép kiểm cho nó ở quy mô bạn
    chạy được, và nói trước kết quả nào sẽ làm bạn nghi ngờ. Chương này đã nêu
    hai chỗ đáng ngờ, bảng A4 và bảng 2, và cả hai đều là chuyện tái lập chứ
    không phải chuyện phương pháp; chỗ phương pháp thì chưa ai trong sách này
    đụng tới.

  \item Chương này không đo được luận điểm của chính nó, và
    mục~\ref{sec:ch09-khong-do-duoc} nói corpus phải đổi chứ không phải script.
    Dựng corpus ấy. Cần một văn phạm sinh ra một miền rộng và một miền hẹp nằm
    trong nó, sao cho miền hẹp có ít dữ liệu và mang cấu trúc mà chỉ miền rộng
    dạy được. Rồi chạy hai pha thật và trả lời: có chuyển giao không, và bạn
    phân biệt \enquote{có chuyển giao} với \enquote{huấn luyện lâu hơn} bằng
    phép đối chứng nào? Bài tập 1 tier ba của
    chương~\ref{ch:mot-vector-cho-moi-tu} đã đòi một văn phạm khó hơn vì một lý
    do khác; nếu bạn dựng nó rồi thì bài này nối tiếp được.

  \item Cả hai bài \tn{quy luật scaling}{scaling law} đều khớp một hàm lũy thừa
    lên mất mát và
    không ai trong hai bài hỏi vì sao mất mát lại là hàm lũy thừa. Đọc
    \eqref{eq:ch09-chinchilla} như một phát biểu về cấu trúc chứ không như một
    phép khớp: số hạng \(E\) nói gì, và hai số hạng còn lại nói gì về chỗ hết
    tham số so với chỗ hết dữ liệu? Rồi đề xuất một dạng hàm khác khớp cùng dữ
    liệu và cho một dự đoán \emph{khác} ngoài dải đã khớp, và nói phép đo nào
    phân biệt được hai dạng.

  \item Mục~\ref{sec:ch09-sau-nd} đo được \(6ND\) sai 44\% ở ngữ cảnh 32768, và
    mục~\ref{sec:ch08-cho-vo} đo được chi phí thật của phần bậc hai đến từ chỗ
    ma trận điểm số rời cache chứ không từ số phép tính. Hai chuyện ấy là hai
    đại lượng khác nhau cùng phụ thuộc \(n\). Câu hỏi: một công thức ngân sách
    cho mô hình ngữ cảnh dài nên đếm cái nào? Đề xuất một thay thế cho \(6ND\)
    trả lời được cả hai, định nghĩa nó đủ chặt để tính được, rồi tính cho GPT-3
    ở bốn độ dài ngữ cảnh và nói công thức của bạn mang thêm tham số phần cứng
    nào.
\end{bt}
